\section{The spectral measure of a single field}
\label{sec:trace-spectrum}
The joint translation measure acts on vectors in Hilbert space. A point
field applied to the vacuum need not be such a vector. To obtain the
spectral measure of its two-point function, we first keep a positive heat
height, take the measure of the resulting vector, and then remove the
exponential damping. The resulting measure can have infinite total mass.

\subsection{Exact damping of vector measures}
We first record a uniqueness fact for bounded holomorphic functions that
will identify heat evolution inside the spectral integral. Suppose $h$ is
bounded and continuous on the closed upper half-plane, holomorphic in its
interior, and zero on the real axis. Fix $x\in\mathbb R$, $s\geq0$ and
apply the Gaussian contour argument from
Lemma~\ref{lem:paley-positive} to $h(z+x)$ with $b=-2\varepsilon s$.
The bottom integral vanishes. On the top edge,
\[
 e^{-ib(t+is)-\varepsilon(t+is)^2}
       =e^{-\varepsilon s^2}e^{-\varepsilon t^2}.
\]
The nonzero constant factor can be cancelled, giving
$\int e^{-\varepsilon t^2}h(x+t+is)dt=0$.
Set $\varepsilon=R^{-2}$ and $t=Ru$. After cancelling $R>0$, let
$R\downarrow0$. Boundedness and the Gaussian majorant justify dominated
convergence and give $\sqrt\pi h(x+is)=0$. Thus $h=0$. In particular,
two such bounded holomorphic functions with equal continuous boundary
values agree throughout the half-plane.

Let $\nu_v$ be the finite joint spectral measure of a vector
$v\in\mathcal H_E$, and retain $T_E$ and $E_t$ from the earlier
construction. Equation~\eqref{eq:euclidean-positive-energy} gives
$p_0\geq0$ almost everywhere for $\nu_v$. For fixed $a\in\mathbb R^4$,
the function
\[
 G(z)=\int e^{ip\cdot a}e^{izp_0}\,d\nu_v(p)
\]
is bounded and continuous on the closed upper half-plane. On a neighborhood
where $\operatorname{Im}z\geq\eta>0$, its derivative is dominated by
$p_0e^{-\eta p_0}\leq\eta^{-1}$. Since $\nu_v$ is finite, dominated
differentiation proves holomorphy. The bounded holomorphic function
\[
 F(z)=\langle v,T_E(a)E(-iz)v\rangle
\]
in the open half-plane has the continuous boundary
$\langle v,T_E(a)V(t)v\rangle$. By the translation spectral formula
these boundary values equal $G(t)$. The uniqueness just proved therefore
identifies $F=G$.

The operators $E_t$ are self-adjoint and commute with translations.
For $t\geq0$, the preceding identity at $z=2it$ gives
\[
 \langle E_tv,T_E(a)E_tv\rangle
  =\langle v,T_E(a)E_{2t}v\rangle
  =\int e^{ip\cdot a}e^{-2tp_0}\,d\nu_v(p).
\]
Both measures here are finite, because $p_0\geq0$. Uniqueness of finite
measures from their characteristic functions, proved in
Section~\ref{sec:spectral-measure}, yields
\begin{equation}\label{eq:trace-vector-damping}
 d\nu_{E_tv}(p)=e^{-2tp_0}\,d\nu_v(p).
\end{equation}
This proof uses the continuous bounded half-plane boundary of a vector
matrix coefficient, not the generally distributional boundary of a point
correlation.

\subsection{Point vectors and removal of damping}
Fix a label $v$ and, for $s>0$, define
\[
 q_s(v)=Q_{p,(s)}(0;v)=E_{s-p}B_p(v)\Omega_E,
 \qquad 0<p<1,\quad4p<s.
\]
The vector is independent of the admissible width $p$. Indeed, the mixed
Gram identity gives the same number for the scalar product of either of
two such vectors with either one of them. Expanding the squared norm of
their difference therefore gives zero. If $t\geq0$, a width admissible
at $s$ is also admissible at $s+t$. Using that common width and the
semigroup law gives
\begin{equation}\label{eq:trace-point-flow}
 q_{s+t}(v)=E_tq_s(v).
\end{equation}
Let $\nu_s$ be the finite spectral measure of $q_s(v)$. Define a positive
Borel measure by
\begin{equation}\label{eq:trace-undamping}
 d\mu_v(p)=e^{2p_0}\,d\nu_1(p).
\end{equation}
This is a well-defined sigma-finite measure: the density is finite and
positive, and is bounded on each set $\{|p_0|\leq n\}$.
Multiplication of densities and
\eqref{eq:trace-vector-damping}--\eqref{eq:trace-point-flow} imply
\[
 e^{2(s+t)p_0}\,d\nu_{s+t}=e^{2sp_0}\,d\nu_s.
\]
Compare $s$ and $1$ through their common larger height $s+1$. It follows
that for every $s>0$,
\begin{equation}\label{eq:trace-height-measure}
 d\mu_v=e^{2sp_0}\,d\nu_s,
 \qquad e^{-2sp_0}\,d\mu_v=d\nu_s.
\end{equation}
In particular $\mu_v$ is supported in the time half-space $p_0\geq0$,
and every positive exponential damping has finite mass.

The single-insertion case of the mixed Gram identity and translation
covariance of the point vectors now gives
\begin{equation}\label{eq:trace-pair-laplace}
 F_2(x_0^0-is,x_1^0+is;\mathbf x_0,\mathbf x_1;Jv,v)
 =\int e^{-2sp_0}e^{ip\cdot(x_1-x_0)}\,d\mu_v(p).
\end{equation}
To see the relative sign explicitly, the two vectors are
$T_E(x_0)q_s(v)$ and $T_E(x_1)q_s(v)$. Unitarity moves the first
translation to the second, producing $T_E(x_1-x_0)$; the spectral formula
and \eqref{eq:trace-height-measure} give the integral above.

For a self-adjoint Euclidean scalar label $v$, taking both real positions
equal to zero identifies the left side with its Euclidean two-point
function at separation $2s$. Thus
\begin{equation}\label{eq:trace-positive-laplace}
 T_v(t)=\int e^{-tp_0}\,d\mu_v(p)\quad(t>0).
\end{equation}
It follows already that $T_v$ is real, nonnegative and decreasing. The
stronger forward-cone conclusion below is obtained after identifying the
measure with the actual reconstructed fields.

\subsection{Polynomial growth and the full boundary distribution}
The fixed-label gap bound \eqref{eq:gap-power}, applied at heights $(-s,s)$
and zero real positions, gives constants $C\geq0$ and $k\in\mathbb N$
such that
\[
 \nu_s(\mathbb R^4)=\|q_s(v)\|^2\leq Cs^{-k},\qquad0<s\leq1.
\]
For $R\geq1$ use $s=R^{-1}$ in \eqref{eq:trace-height-measure}. On
$\{p:p_0\leq R\}$ one has $e^{-2sp_0}\geq e^{-2}$, and hence
\begin{equation}\label{eq:trace-energy-growth}
 \mu_v\{p:p_0\leq R\}\leq e^2CR^k.
\end{equation}
This controls an entire energy sublevel set, not just a Euclidean ball.
In particular $\mu_v(\{|p|\leq R\})\leq AR^k$ with $A=e^2C$.
The measure is locally finite and therefore is a Radon measure on
$\mathbb R^4$.

Here is a quantitative rapid-weight consequence. Cover $\mathbb R^4$ by
the closed shells $2^j\leq1+|p|\leq2^{j+1}$, $j\geq0$. On the $j$th
shell the weight $(1+|p|)^{-k-1}$ is at most $2^{-j(k+1)}$, and its
mass is at most $A2^{(j+1)k}$. Summing, with overlaps causing only an
upper bound, gives
\begin{equation}\label{eq:trace-rapid-weight}
 \int(1+|p|)^{-k-1}\,d\mu_v(p)\leq A2^{k+1}<\infty.
\end{equation}
Thus every Schwartz function is absolutely integrable against $\mu_v$,
and integration is a continuous functional on Schwartz space, bounded
by a weighted supremum seminorm times the finite integral in
\eqref{eq:trace-rapid-weight}.

Use the unnormalized plus Fourier convention
\[
 f_+(p)=\int e^{ip\cdot x}f(x)\,dx,
 \qquad
 \mathcal W\Phi(p)=\int e^{ip\cdot(x_1-x_0)}\Phi(x_0,x_1)\,dx_0dx_1.
\]
Both are Schwartz functions when their arguments are Schwartz.
For the second assertion take the joint plus Fourier transform and
restrict it to the injective linear subspace $p\mapsto(-p,p)$.
The chain rule controls derivatives of the restriction, and
$|(-p,p)|=\sqrt2|p|$ controls its weighted supremum seminorms.

At each height $s>0$, Fubini in \eqref{eq:trace-pair-laplace} is justified
by $\nu_s(\mathbb R^4)\|\Phi\|_{L^1}<\infty$. It identifies the joint
tested correlation with
$\int e^{-2sp_0}\mathcal W\Phi(p)d\mu_v(p)$.
As $s\downarrow0$, the exponential has absolute value at most one
almost everywhere and tends to one. The integrable majorant
$|\mathcal W\Phi|$ and the joint boundary theorem give
\begin{equation}\label{eq:trace-full-boundary}
 W_2(Jv,v)(\Phi)=\int\mathcal W\Phi(p)\,d\mu_v(p).
\end{equation}
This is an identity on all joint Schwartz tests, including those meeting
the coincidence set or the null cone. No restriction of a distribution
to such a set has been used.

\subsection{Identification with field vectors and the forward cone}
Let $\xi_v(f)=\phi(v,f)\Omega$ be the one-field vector in the reconstructed
Hilbert space. With
$\Phi(x_0,x_1)=\overline{f(x_0)}g(x_1)$, Fubini in the ordinary Fourier
integral gives
$\mathcal W\Phi=\overline{f_+}g_+$. The word Gram formula and
\eqref{eq:trace-full-boundary} therefore give
\begin{equation}\label{eq:trace-smeared-inner}
 \langle\xi_v(f),\xi_v(g)\rangle
                  =\int\overline{f_+(p)}g_+(p)\,d\mu_v(p).
\end{equation}
The product is integrable by the rapid-weight bound.
Translation covariance gives
$U(a)\xi_v(f)=\xi_v(f(\,\cdot-a))$, and changing variables in the
Fourier integral gives
$(f(\,\cdot-a))_+(p)=e^{ip\cdot a}f_+(p)$.
It follows from \eqref{eq:trace-smeared-inner} that the characteristic
function of the finite PVM measure of $\xi_v(f)$ equals that of
$|f_+|^2\mu_v$. The latter is finite by
\eqref{eq:trace-rapid-weight}. Finite-measure uniqueness gives
\begin{equation}\label{eq:trace-field-mass}
 d\mu^{\mathrm{field}}_{\xi_v(f)}=|f_+|^2\,d\mu_v.
\end{equation}

The weights in this identity determine $\mu_v$ uniquely. Given a closed
ball of radius $n$, choose a smooth compactly supported function $h$ equal
to one on that ball and take a Schwartz function $f$ with $f_+=h$.
Such an $f$ exists by Fourier inversion, with the factor $(2\pi)^{-4}$
in the inverse transform. The weight $|f_+|^2$ is then exactly one on the
ball. Equality of all the weighted measures implies equality on every
measurable subset of each ball; increasing the radius proves equality
on all Borel sets by continuity from below.

By Theorem~\ref{thm:field-spectrum}, every measure on the left of
\eqref{eq:trace-field-mass} vanishes outside the forward cone.
Apply the preceding weights on each ball to conclude
\begin{equation}\label{eq:trace-forward}
 \mu_v(\mathbb R^4\setminus\overline V_+)=0.
\end{equation}
Thus the forward-cone statement is inherited from the reconstructed field
representation, while the initial undamping argument required only
positivity of the time coordinate in the Euclidean chamber.

For completeness the same identification gives Lorentz covariance of the
point measure. Write $A=\Lambda^{-T}$ for the dual momentum action and
$g_\Lambda(x)=g(\Lambda^{-1}x)$. Since $\det\Lambda=1$,
$(g_\Lambda)_+(Ap)=g_+(p)$. Covariance of the field vectors and of their
PVM measures then gives
\[
 |(g_\Lambda)_+|^2\mu_{\rho(\Lambda)v}
            =A_*\bigl(|g_+|^2\mu_v\bigr)
            =|(g_\Lambda)_+|^2 A_*\mu_v.
\]
The transformed tests run through Schwartz space, so uniqueness by the
weights proves $\mu_{\rho(\Lambda)v}=A_*\mu_v$. In particular a Lorentz
scalar label has a Lorentz-invariant point measure, with the same dual
convention as the joint translation measure.

\subsection{Application to the centered trace}
Let $T_{ab}$ be the real symmetric Euclidean stress tensor supplied by the
construction, transforming as a rank-two tensor, and put
$\tau=\sum_aT_{aa}$. Orthogonality gives
\[
 R(K)\tau=\sum_{a,i,j}K_{ia}K_{ja}T_{ij}
          =\sum_{i,j}\delta_{ij}T_{ij}=\tau.
\]
Its one-point value $c=\langle\tau\rangle_\chi$ is independent of position
by translation invariance, and real by conjugation symmetry. The centered
label $\Theta=\tau-c\mathbf1$ is consequently real, invariant under all
Euclidean orthogonal transformations, and has zero one-point function.
Since the adjoint label operation is conjugation followed by time
reflection, $J\Theta=\Theta$.

It is also a Lorentz scalar. Differentiating its Euclidean invariance shows
that every Euclidean rotation generator annihilates it. The corresponding
boost generator is $-i$ times that generator, so also annihilates it. The
constant curve $\Theta$ and its boost orbit solve the same finite-dimensional
constant-coefficient differential equation with the same initial value;
uniqueness makes the orbit constant. Spatial rotations already fix it,
and generation by boosts and spatial rotations proves full Lorentz
invariance.

The one-point holomorphic continuation of $\Theta$ is zero: on positive
imaginary times it is the zero Euclidean one-point function, and the
one-variable identity theorem extends that equality to the half-plane
for each spatial position. Its tested boundary is therefore zero. The
word formula gives
$\langle\Omega,\phi(\Theta,f)\Omega\rangle=0$ for every Schwartz $f$.
Finally \eqref{eq:trace-positive-laplace} supplies the positive measure
and Laplace formula for the very same centered trace used in
Section~\ref{sec:trace-asymptotic}. No additional positive-measure input
is needed there once the earlier reconstruction hypotheses and this
stress tensor are available.
